% How the proof of Theorem 1.1 is put together.
\begin{tikzpicture}[
    box/.style={draw=black!70, line width=0.6pt, rounded corners=3pt, fill=#1,
      text width=10.3cm, align=flush left, inner sep=5pt, font=\small},
    box/.default=white,
    arr/.style={-{Stealth[length=6pt]}, line width=0.8pt}]
  \node[box=gray!8] (cube) {\textbf{Lemma~\ref{lem:hjcube}.} If $|Z|\ge M(t)$, every $2$-colouring of
    the subsets of $Z$ has a monochromatic $t$-cube $\{U\cup X_I: I\subseteq[t]\}$ inside $Z$.};
  \node[box=teal!8, below=8mm of cube] (sat) {\textbf{Proposition~\ref{prop:saturation}.} For every
    $r\le n-M$ some shape $(u;a_1,\dots,a_t)$ with $r\le u\le r+M-t$ satisfies
    $\rho_\chi(u;a)\ge(t+2)^{-M}$.};
  \node[box=teal!13, below=8mm of sat] (slice) {\textbf{Lemma~\ref{lem:slices}.} For each level $j$
    and generator set $A$, at least $M^{m-t-1}$ of the $M^m$ size vectors $\beta$ give the cube
    $Q(j,A)$ a good shape.};
  \node[box=teal!18, below=8mm of slice] (avg) {\textbf{Lemma~\ref{lem:average}.} Some ordering
    $\sigma$ and sizes $\beta$ make at least \mbox{$|\mathcal P|/((t+2)^M M^{t+1})$} of the cubes
    monochromatic at once.};
  \node[box=teal!24, below=8mm of avg] (glue) {\textbf{Proposition~\ref{prop:glue}.} Monochromatic
    cubes of one block chain glue into a monochromatic union-closed family, so
    \mbox{$F(n)\ge |\mathcal P|/(2(t+2)^M M^{t+1})\ge c_t\,n^{t+1}$}.};
  \foreach \a/\b/\t in {cube/sat/{one of $(t+2)^M$ patterns serves many $\sigma$},
      sat/slice/{tune $t+1$ block sizes, fix the rest},
      slice/avg/{average over $(\sigma,\beta)$},
      avg/glue/{union rule, Lemma~\ref{lem:union}}}
    \draw[arr] (\a) -- node[right=3pt, font=\footnotesize\itshape, text=black!70] {\t} (\b);
\end{tikzpicture}
